% ID: SEL-D0020-OA-E03
% ESTADO: cotejada. Segunda pasada visual completa y borrador compilado revisados.
% Fuente: Selectividad/matematicas_aplicadas/reserva b/Examen-Matemáticas Aplicadas a las Ciencias Sociales.pdf
% SHA-256: 37584a9604930b74008a6bda88ef3d61fa926d5e6a540ab9013569258f772256
% Página: 1; opción A; ejercicio original 3.

En un polideportivo municipal hay inscritos 520 usuarios de los que 220 son niños, 208 son adultos menores de 60 años y el resto adultos mayores de 60 años. De los inscritos, $1/5$ de los niños, el 75\% de los adultos menores de 60 años y 23 adultos mayores de 60 años utilizan las duchas normalmente. Se elige un usuario al azar.
\begin{enumerate}[label=\alph*)]
\item \textbf{(1 punto)} Calcule la probabilidad de que se duche en las instalaciones del polideportivo.
\item \textbf{(0.75 puntos)} Calcule la probabilidad de que sea adulto menor de 60 años y utilice las duchas.
\item \textbf{(0.75 puntos)} Sabiendo que utiliza las duchas, halle la probabilidad de que sea un niño.
\end{enumerate}
